\section{Un flux certain}
\label{sec:U1}

Toute la comptabilité prudentielle repose sur une seule opération, répétée et
généralisée : ramener à aujourd'hui la valeur d'un montant disponible plus tard.
Cette unité en pose la forme la plus simple. Les cinq suivantes ne feront
qu'y ajouter, une par une, les complications du monde réel.

\subsection{La notion : valeur actuelle d'un flux unique}

\begin{notion}{Actualisation et capitalisation}
Un montant $C$ disponible dans $n$ années n'a pas la même valeur qu'un montant
$C$ disponible aujourd'hui. Au taux annuel $r$, un euro placé aujourd'hui vaut
$(1+r)^n$ euros dans $n$ ans : c'est la \emph{capitalisation}\index{capitalisation}.

L'opération inverse, l'\emph{actualisation}\index{actualisation}, ramène un flux
futur à sa valeur d'aujourd'hui :
\[
  V_0 \;=\; \frac{C}{(1+r)^n} \;=\; C \times v^{\,n},
  \qquad\text{avec}\qquad v \;=\; \frac{1}{1+r}.
\]
Le terme $v^{\,n}$ est le \emph{facteur d'actualisation}\index{facteur d'actualisation}
à l'horizon $n$. Il est toujours compris entre $0$ et $1$ dès que $r>0$.
\end{notion}

Deux propriétés doivent devenir des réflexes, car elles se retrouveront à chaque
étage de la progression : la valeur actuelle \emph{décroît} quand le taux
augmente, et elle \emph{décroît} quand l'échéance s'éloigne.

\begin{center}
\begin{tikzpicture}[x=1.35cm, y=1cm]
  \draw[timeaxis] (0,0) -- (5.4,0);
  \foreach \x/\l in {0/0, 1/1, 2/2, 3/{\dots}, 4/n}
    { \node[timenode] at (\x,0) {}; \node[etiq, below=3pt] at (\x,0) {\l}; }
  \draw[flux] (4,0) -- (4,1.15);
  \node[etiq, above] at (4,1.15) {$C$};
  \draw[fluxneg] (0,0.95) -- (0,0.05);
  \node[etiq, above] at (0,0.95) {$V_0 = C\,v^{\,n}$};
  \draw[->, dashed, grisactu] (3.85,1.0) to[bend right=18] (0.15,1.0);
\end{tikzpicture}
\end{center}

\begin{exercice}[valeur actuelle]
Un flux certain de $25\,000$~\euro{} sera perçu dans $7$~ans. Le taux
d'actualisation annuel est de $3{,}20\,\%$. Calculer sa valeur actuelle.
\end{exercice}

\begin{corrige}
\[
  V_0 = \frac{25\,000}{(1{,}032)^{7}} = \frac{25\,000}{1{,}24669} \approx 20\,053\ \text{\euro}.
\]
Autrement dit, disposer de $20\,053$~\euro{} aujourd'hui équivaut à disposer de
$25\,000$~\euro{} dans sept ans, si l'on peut placer son argent à $3{,}20\,\%$.
\end{corrige}

\begin{exercice}[capitalisation]
Un capital de $15\,000$~\euro{} est placé pendant $12$~ans au taux annuel de
$2{,}50\,\%$. Quel montant sera disponible au terme ?
\end{exercice}

\begin{corrige}
\[
  C_{12} = 15\,000 \times (1{,}025)^{12} = 15\,000 \times 1{,}34489 \approx 20\,173\ \text{\euro}.
\]
C'est l'opération réciproque de la précédente : on multiplie au lieu de diviser.
\end{corrige}

\subsection{Généralisation : les conventions de taux}

\begin{generalisation}{Un taux n'a de sens qu'avec sa convention}
Annoncer « $4\,\%$ » ne suffit pas. Il faut préciser la périodicité de
capitalisation et la règle de comptage des jours. Deux distinctions suffisent
pour la suite du cours.

\medskip
\textbf{Taux proportionnel et taux équivalent.} Pour passer d'un taux annuel
$r$ à un taux sur une fraction $1/k$ d'année, deux conventions coexistent :
\[
  \underbrace{r^{\text{prop}}_{k} = \frac{r}{k}}_{\text{proportionnel}}
  \qquad\text{et}\qquad
  \underbrace{r^{\text{éq}}_{k} = (1+r)^{1/k} - 1}_{\text{équivalent}}.
\]
Seul le taux équivalent respecte la capitalisation : $(1+r^{\text{éq}}_{k})^{k} = 1+r$.
Le taux proportionnel, lui, surestime toujours le rendement réel.

\medskip
\textbf{Base de calcul des jours.} /!\ Attention /!\ La convention des marché monétaire est que pour les durées inférieures à un an, on n’utilise pas l’actualisation pour le calcul des intérêts courus, mais la formule suivante : 
\[
  I = C \times r \times \frac{\text{nombre de jours}}{\text{base}}.
\]
La base de calcul des jours peut être de $360$ ou $365$ jours : \emph{exact/360}\index{exact/360} (base monétaire) ou \emph{exact/365}\index{exact/365} (base obligataire) ; la base $360$ produit mécaniquement des intérêts plus élevés que la base $365$ pour un même taux affiché.
\end{generalisation}

\begin{exercice}[proportionnel contre équivalent]
Le taux annuel est de $4{,}00\,\%$. Calculer le taux semestriel proportionnel et
le taux semestriel équivalent. Un capital de $200\,000$~\euro{} est placé un an,
avec capitalisation semestrielle : chiffrer l'écart entre les deux conventions.
\end{exercice}

\begin{corrige}
Taux proportionnel : $r^{\text{prop}}_{2} = 4{,}00/2 = 2{,}0000\,\%$.

Taux équivalent : $r^{\text{éq}}_{2} = (1{,}04)^{1/2} - 1 = 1{,}98039\,\%$.

\medskip
Au bout d'un an, soit deux semestres :
\begin{align*}
  \text{convention proportionnelle} &: 200\,000 \times (1{,}02)^{2} = 208\,080{,}00\ \text{\euro} \\
  \text{convention équivalente}     &: 200\,000 \times (1{,}04)\phantom{^{2}} = 208\,000{,}00\ \text{\euro}
\end{align*}
L'écart est de $80$~\euro{} sur un an. Il paraît anodin ; rapporté à un
portefeuille de provisions techniques de plusieurs centaines de millions et à un
horizon de projection de trente ans, il ne l'est plus.
\end{corrige}

\begin{exercice}[bases de calcul]
Un placement de $500\,000$~\euro{} est effectué pour $91$~jours au taux de
$3{,}00\,\%$. Calculer les intérêts en base exact/360 puis en base exact/365, et
commenter l'écart.
\end{exercice}

\begin{corrige}
\begin{align*}
  \text{exact}/360 &: 500\,000 \times 0{,}03 \times \tfrac{91}{360} = 3\,791{,}67\ \text{\euro} \\
  \text{exact}/365 &: 500\,000 \times 0{,}03 \times \tfrac{91}{365} = 3\,739{,}73\ \text{\euro}
\end{align*}
L'écart atteint $51{,}94$~\euro, soit $1{,}39\,\%$ des intérêts. Le rapport des
deux bases vaut exactement $365/360$, quel que soit le montant, le taux ou la
durée : c'est un écart de convention, pas un écart économique.
\end{corrige}

\begin{pont}{Où cela sert dans Solvabilité~II}
Le facteur $v^{\,n}$ de cette unité est très exactement celui qui sera appliqué
à chaque flux de prestation dans le calcul de la meilleure estimation. La seule
différence, à ce stade purement technique, est que le taux $r$ n'y est pas choisi
librement : il est lu sur la courbe des taux sans risque publiée chaque mois par
l'EIOPA, et il varie selon l'échéance du flux. C'est l'objet de l'unité~3.
\end{pont}

\begin{refreg}
\textbf{Références.} Directive 2009/138/CE, article~77, paragraphe~2 : la
meilleure estimation est la moyenne pondérée des flux de trésorerie futurs
compte tenu de la valeur temporelle de l'argent, obtenue au moyen de la courbe
des taux sans risque pertinente. \\
Les structures par terme sans risque sont publiées mensuellement par l'EIOPA,
accompagnées de leur documentation technique.
\end{refreg}

\subsection*{Notations introduites}
\begin{description}[leftmargin=3cm, style=nextline, font=\normalfont]
  \notation{C}{montant nominal d'un flux}
  \notation{r}{taux d'actualisation annuel}
  \notation{n}{horizon, en années}
  \notation{v = (1+r)^{-1}}{facteur d'actualisation à un an}
  \notation{V_0}{valeur actuelle}
\end{description}
